% =====================================================================
%  Section V. Weak-Coupling Approximation
% =====================================================================
\section{WEAK-COUPLING APPROXIMATION}\label{sec:weak}

The weak-coupling approximation will be discussed briefly, leaving many
questions open.  Only the pure gauge field will be discussed.  Consider
again the expression
\begin{equation}
Z_{\theta}(P) = \prod_{n,\mu\nu}\int\ddth{\mu\nu}
  \exp\!\Bigl[i\sum_{P}(\pm)\thl{\mu\nu}
    +\frac{1}{2g^{2}}\sum_{n,\mu\nu}
      \bigl(1-\tfrac{1}{2}\bigl(e^{if_{\mu\nu}(n)}+e^{-if_{\mu\nu}(n)}\bigr)\bigr)\Bigr].
\label{eq:5.1}
\end{equation}

Suppose the integration variables were $f_{\mu\nu}$ rather than
$\thl{\mu\nu}$.  For small $g$, only small values of $f_{\mu\nu}$ would
be important in the integral, in order that
$\exp[(1/2g^{2})\cos f_{\mu\nu}]$ be near its maximum value 1.  One
would then expand:
\begin{equation}
\frac{1}{2g^{2}}\bigl(1-\cos f_{\mu\nu}\bigr)
  \approx \frac{1}{4g^{2}}f_{\mu\nu}^{2}+\cdots .
\label{eq:5.2}
\end{equation}
With this approximation one could extend the limits of integration on
$f_{\mu\nu}$ from $\pm\pi$ to $\pm\infty$, with negligible error; one
would then have a set of Gaussian integrals to evaluate.

In practice the integration variables are the $\thl{\mu\nu}$, not the
$f_{\mu\nu}$.  However, one can make a change of variable from the
$\thl{\mu\nu}$ to the $f_{\mu\nu}$.  It is not possible to eliminate all
the $\thl{\mu\nu}$ by this transformation, and not all the variables
$f_{\mu\nu}$ are independent.  Nevertheless, the transformation is
sufficient to make $I_{\theta}(P)$ calculable for small $g^{2}$.

To make the change of variables precise, consider a system of finite
size ($|n_{0}|,|n_{1}|,|n_{2}|,|n_{3}| \le N$) with periodic boundary
conditions.  Then one can change variables from the $\thl{\mu\nu}$ to a
subset of the $f_{\mu\nu}$ plus some gauge transformation variables
$\varphi_{n}$, plus four extra variables $\xi_{\mu\nu}$, as follows:

\emph{(i)}  For $n$ with $n_{0} < N$, $\thl{\mu\nu}$
($\mu=1,2,3$) is replaced by $f_{\mu\nu}$.  For $\thl{0\mu}$, one writes
\begin{equation}
\thl{0\mu} \to f_{0\mu},
\label{eq:5.3}
\end{equation}
and replaces $\thl{0\mu}$ by $\varphi_{\mu}$.  This is the essence of the
transformation from $\thl{\mu\nu}$ to $f_{\mu\nu}$ for $\mu \ne 0$ and
from $\thl{0\mu}$ to $\varphi_{\mu}$.  To complete the transformation one
must discuss the surface $n_{0}=N$.

\emph{(ii)}  For $n_{0}=N$, $n_{1}<N$, and $n_{2}$ and $n_{3}$ arbitrary,
$\thl{\mu\nu}(\mu,\nu\in\{1,2,3\})$ is replaced by $f_{\mu\nu}$;
$\thl{0\mu}$ is replaced by $\varphi_{\mu}$, with
\begin{equation}
\thl{0\mu}(n) = \varphi_{\mu}(n) + i\,\text{(linear combination of $f_{\mu\nu}$)}.
\label{eq:5.4}
\end{equation}

\emph{(iii)}  For $n_{0}=n_{1}=N$, $n_{2}<N$, and $n_{3}$ arbitrary,
$\thl{\mu\nu}(\mu,\nu\in\{2,3\})$ is replaced by $f_{\mu\nu}$;
$\thl{0\mu}$ is replaced by $\varphi_{\mu}$, with
$\thl{0\mu}=f_{0\mu}-\varphi_{\mu}$.

\emph{(iv)}  For $n_{0}=n_{1}=n_{2}=N$, and $n_{3}<N$,
$\thl{\mu\nu}(\mu=3)$ is replaced by $f_{3\nu}$; $\thl{0\mu}$ by
$\varphi_{\mu}$, where $\thl{0\mu}=\varphi_{\mu}-\varphi_{\mu}$.

\emph{(v)}  For $n_{0}=n_{1}=n_{2}=n_{3}=N$ one writes
$\thl{\mu\nu}=\xi_{\mu\nu}$ with $\xi_{\mu\nu}$ being the new variables.
One also sets $\varphi_{n}=0$ for this value of $n$.

The variables $f_{\mu\nu}$ which are not integration variables (for
example, $f_{\mu\nu}$ with $\mu\ne0$ and $\nu\ne0$, for $n$ with
$n_{0}<N$) can be expressed in terms of the independent $f_{\mu\nu}$
variables using Eq.~\eqref{eq:4.6}; neither the $\varphi_{n}$ nor
$\xi_{\mu\nu}$ appear in these expressions.

When an arbitrary $\thl{\mu\nu}$ is expressed in terms of the new
independent variables, one finds ($\xi_{\mu\nu}$ is present only if
$n_{0}=N$):
\begin{equation}
\begin{split}
\thl{\mu\nu}(n) = f_{\mu\nu}(n) &+ \text{(gauge-transformation terms)}\\
 &+ \text{(linear combination of $f_{\mu\nu}$)},
\end{split}
\label{eq:5.5}
\end{equation}
i.e., the $\varphi$ variables define a gauge transformation and $f$
represents a translation of some of the $\theta$'s.  It is easily
verified that the integrand of $Z_{\theta}(P)$ involves only the $f$'s:
It is independent of both the $\varphi$'s and $\xi_{\mu\nu}$ (the latter
does not appear because any closed path $P$ has as many
$-\thl{\mu\nu}$ terms as $+\thl{\mu\nu}$ terms on the sublattice
$n_{0}=N$).  Hence the $\varphi_{n}$ and $\xi_{\mu\nu}$ integrations can
be computed trivially.  The integrations are nontrivial because of the
constraint~\eqref{eq:4.6}.

What one wants to accomplish is to reduce the lattice theory for small
$g$ to something like a conventional free gauge-field theory.  This
means restoring the $\thl{\mu\nu}$ as the integration variables, but
with infinite limits of integration, and with a gauge-fixing term
included.  Suppose, for example, one starts with
\begin{equation}
\begin{split}
Z_{\theta}'(P) = \prod_{n,\mu\nu}\int\ddth{\mu\nu}\,
  \exp\!\Bigl[i\sum_{P}(\pm)\thl{\mu\nu}
    &+\frac{1}{2g^{2}}\sum_{n,\mu\nu}
      \bigl(1-\tfrac{1}{2}\bigl(e^{if_{\mu\nu}(n)}+e^{-if_{\mu\nu}(n)}\bigr)\bigr)\\
    &-\frac{\alpha}{2}\sum_{n}\bigl(\nabla_{\mu}A_{\mu}(n)\bigr)^{2}\Bigr],
\end{split}
\label{eq:5.6}
\end{equation}
where the $\alpha$ term is a lattice version of a
$(\nabla_{\mu}A_{\mu})^{2}$ gauge-fixing term.  This integral can be
computed by explicit Gaussian integration methods rather more easily
than the $f_{\mu\nu}$ integrations for $I_{\theta}(P)$.  In addition,
$I_{\theta}(P)$ can also be reduced to an integral over a subset of the
$f_{\mu\nu}$, using the same change of variables as for $I_{\theta}(P)$.
The result is different in this case due to the $\alpha$ term which
couples the $\varphi$'s to the $f$'s; also there is no convergence
factor for the $\xi$ integration.  To make $I_{\theta}'(P)$ well defined
and equal to $I_{\theta}(P)$, one must (a) put in a convergence factor
for the $\xi$ integral, i.e., a term $-\tfrac{1}{2}\alpha(\nabla_{\mu}A_{\mu})^{2}$, and
(b) add a quadratic form in the $\xi$'s to compensate for the result of
the $\xi$ integration of the gauge-fixing term.  The author has not
carried through this calculation; but since the net result is still that
$Z_{\theta}(P) = I_{\theta}'(P)$ is a Gaussian integration in the
$\theta$'s, the result will presumably be similar to the conventional
free-field calculation reported in Sec.~\ref{sec:binding}.
